\section{Pre-training：OLMo 2 的可复现经验}

\subsection{目标：在可控成本内做出高竞争力底座}
讲者用 OLMo 2 讨论了一个关键现实：研究型团队并不总有 frontier 级算力，但依然可以通过更好的数据治理与训练细节，在 7B/13B 规模做出强基线。

\begin{importantbox}{OLMo 2 的价值在于 ``可解释改进''}
相较只公布结果曲线，OLMo 系列更强调公开训练链路，使社区可以系统回答：
\begin{itemize}
    \item 哪些训练细节真正贡献了性能；
    \item 哪些改动只在特定基准有效；
    \item 在不同预算下哪组配置更稳健。
\end{itemize}
\end{importantbox}

\begin{figure}[H]
\centering
\includegraphics[width=0.9\textwidth]{frame-03-olmo-scaling.jpg}
\caption{视频约 00:18:30：OLMo 训练曲线与规模比较，强调在预算受限下的效率优化}
\end{figure}

\subsection{预训练里真正影响推理底座的因素}
\begin{knowledgebox}{比 ``参数规模'' 更容易被忽略的四类变量}
\begin{enumerate}
    \item 数据去重、污染控制与质量分层；
    \item 领域采样比例（代码、数学、知识文本、对话）；
    \item 学习率与 warmup/cooldown 调度；
    \item tokenizer 与上下文窗口策略。
\end{enumerate}
\end{knowledgebox}

讲者强调，推理能力在预训练阶段已经被部分决定，后训练只能在已有能力边界内做重排和强化。这解释了为何很多后训练技巧在不同底座上效果差异很大。

\begin{warningbox}{预训练指标与下游推理能力并非一一对应}
较低的语言建模 loss 不必然带来更强链式推理。若训练语料中的推理结构信号不足，模型可能擅长流畅续写却不擅长多步约束推断。
\end{warningbox}

\subsection{本章小结}
OLMo 2 的启示是：在研究场景中，公开且可解释的预训练配方本身就是成果，它决定了后续所有 reasoning 增强的上限与可重复性。


